\chapter{Introduction}

JavaScript Object Notation (JSON) is a lightweight, text-based data interchange format widely used for communication between web clients and servers. Its simplicity, readability, and ability of parsing have made it very impotent in today's software systems for exchanging structured data.

Parsing JSON data efficiently and accurately is critical to ensure seamless data transfer and integration between heterogeneous systems. A JSON parser reads JSON text, check its correct or not according to the JSON syntax, and converts it into valid data structures for a program. This process is being an essential part of many applications, web APIs, configuration management, and real-time data processing.

The key goal behind this seminar report is to research about the design and implementation of a JSON parser, understand the parsing techniques, and evaluate parser performance. By developing a custom JSON parser, this report aims to provide knowledge into compiler construction concepts such as lexical analysis, syntax parsing, and error handling relevant to a practical and widely used data format.

The objectives of this seminar report are as follows:
\begin{itemize}
	\item To study the JSON format, its syntax, and grammar rules.
	\item To analyze different parsing techniques applicable for JSON, including recursive descent and LR parsing.
	\item To implement a JSON parser using tools like Lex and Yacc, focusing on accurate parsing and error detection.
\end{itemize}

This report based on the syntax parsing aspects of JSON, error handling mechanisms, and practical implementation details. 

The report is organized as following: Chapter 2 have details existing JSON parsers and parsing techniques. Chapter 4 is about the methodology used for the parser design and implementation. Chapter 5 gives the results, testing, and analysis of the parser. Finally, Chapter 5 is the conclusion  the report and suggests directions for future enhancements.

With this seminar report, readers can understand a comprehensive understanding of JSON parsing and get practical knowledge to compiler construction techniques applied to a real-world problem.
